योगश्चित्तवृत्तिनिरोध इत्युक्तम। तत्र बोद्धृत्वेनैव पुरुषसद्भावाधिगमः। विषयभुतवृत्तिनिरोधे च विषयिणोऽपि बोद्धुः पुरुषस्य निरोधं कश्चिदा\vrtms{शंकेतापि}{}{}। तथा च सति तत्कैवल्यप्राप्त्युपायस्यापि विवेकख्यातेरनर्त्थकत्वम् मन्वीत। योगानुशा\vrtms{सनस्य च}{}{} तदर्त्थस्य निष्फलत्व\vrtms{मपि}{}{} प्रतिपद्येत॥

तस्मात्पुरुषस्य वृत्तिनिरोदहादनिरोधन्दर्शयितुकामः ख्यातेश्चापि फलं साक्षाद्दर्शयन्नाह---\bht{तदवस्थे चेतसी}त्यादि॥

तदवस्थे---निरोधावस्थे। \bht{स्वविषयाभावात्}---स्वविषयश्चित्तवृत्तिः। तदभावे \bht{बुद्धिबोधात्मा}---बुद्धिं वृत्तिरूपेण परिणताम्बुध्यत इति बुद्धेर्बोद्धा---\bht{पुरुषः}। तद्बोधनमेव हि पुरुषस्य स्वरूपम्। नान्यो बोद्धा, नान्यद्बोधनम्। बोद्धुश्च बोधनादन्यत्वे सति विक्रियात्मकता स्यात्। ततश्च दर्शितविषयत्वञ्च न स्यात्। तथा च करणान्तरापेक्षित्वम्पुरुषस्य प्रसज्येत। तस्माद्वोधनम्बोद्धृत्वं चोपचरितम्। त्दवृत्तिसारूप्येन। तथा च वक्षयति---द्रष्टा दृशिमात्र इति। तस्मादाह बुद्धिबोधनमेवात्मा स्वरूपं यस्य स बुद्धिबोधात्मा पुरुषः॥

किंस्वभावः किन्निस्वभावः---सद्भावे वा किंस्वभावः कथं वासद्भाव इति---

\str{तदा द्रष्टुः स्वरूपे ऽवस्थानम्॥३॥}

यदा निरुद्धा वृत्तयः॥
\bht{स्वरूपप्रतिष्ठा तदानीञ्चितिशक्तिः}---\vrtsm{स्वरूपप्रतिष्ठैवे}{\em}{}त्यर्त्थः। \bht{यथा कैवल्ये}। \vrt{कैवल्ये}{}{} सद्भावन्तु `न तत्स्वाभासन्दृश्यत्वादि'त्येवमादिना वक्ष्यति। तत्र सिद्धे सद्भावे तत अकृष्य प्राप्ताशंकानिवृत्यर्त्थमुच्यते---`स्वरूपप्रतिष्ठे'ति॥

सिद्धवद्व्याख्यातं सूत्रम। तत्रेदम्प्रसक्तम्---इतरत्र न स्वरूपप्रतिष्ठेति। अन्यथा हि तदेति विशेषणानर्त्थक्यं स्यात्। न चेदन्यत्र स्वरूपप्रतिष्ठा तत्रावस्थान्तरयोगात्परिणामवत्वादिदोषः प्राप्नोतीति तत्परिजिहीर्षयाहा---\bht{व्युत्थानचित्ते तु सति तथापि भवन्ती}ति। अत्रापि स्वरूपप्रतिष्ठत्वन्दर्शयति। \bht{न तथे}त्यनेन `तदे'ति विशेषणस्यार्त्थवत्वं ख्यापयति॥३॥
